\section{模型假设}

\subsection{问题一的模型假设}

为在5个稀疏观测点下保持模型可识别和结论可解释，作如下假设：

\begin{enumerate}[label=(\arabic*),leftmargin=2em]
  \item 五个已核实数据点使用一致的成人超重肥胖合计口径，能够反映我国总体变化趋势。
  \item 成人超重肥胖率位于$0$--$100\%$之间，主分析将$K=100\%$作为数学物理上限，而不把它解释为未来一定达到的实际平台值。
  \item 研究期内历史增长机制不存在无法由模型结构与上限情景表达的突变；模型只预测全国成人超重肥胖合计比例。
\end{enumerate}

\subsection{问题二的模型假设}

为统一不同研究的证据口径并保证排序可解释性，作如下假设：

\begin{enumerate}[label=(\arabic*),leftmargin=2em]
  \item 系统综述纳入的研究能够代表非手术、非药物成人减重维持的一般证据，用于限定模型的适用人群和干预范围。
  \item 研究数量、统一方向后一致率和Strong/Moderate证据编码均可作为效益型指标，数值越大表示证据优先级越高。
  \item 补充证据因结局或口径不同只用于机制解释；TOPSIS得分不等同于因果效应或个体失败概率。
\end{enumerate}

\subsection{问题三的模型假设}

为保证画像、方案和排序结果具有统一边界，作如下假设：

\begin{enumerate}[label=(\arabic*),leftmargin=2em]
  \item 五类生活方式画像用于表示典型工作生活情境，不等同于对真实个体的医学诊断，也不假设其来自统计聚类。
  \item 四项指南范围只用于候选方案初筛，不能替代考虑疾病、妊娠、用药、性别和基础能量需要量的个体安全评估。
  \item 离散画像、覆盖赋值、90\,kg基准和能量折算均用于组内相对比较；TOPSIS得分不跨画像比较，也不表示成功率。
\end{enumerate}

\subsection{问题四的模型假设}

为区分人口预测、模型结果与政策情景，作如下假设：

\begin{enumerate}[label=(\arabic*),leftmargin=2em]
  \item 联合国WPP 2024中国中方案的单岁年中人口能够用于描述2024---2050年18岁及以上人口规模；其属于人口预测，不视为未来真实观测。
  \item 问题一比例与18岁以上人口年龄口径近似一致；覆盖、条件依从和效果位于$[0,1]$且在模拟中暂视为相互独立。
  \item 低中高参数、模型等权和$K$分布均为情景设定；可避免人数只表示单年现患负担差，不等同于治愈人数或经济损失。
\end{enumerate}

% ==================== 四、符号说明 ====================
